\section{Conclusion}
\label{sec:conclusion}

We introduced \method{}, a sparse-attention mechanism co-designed with Grouped-Query Attention. The architecture attaches a lightweight Index Branch to a standard GQA layer: each GQA group independently selects a small set of key-value blocks through a block-level dot-product indexer, and the Main Branch performs softmax attention restricted to the selected blocks. The Index Branch is a pure selector and is trained by a KL alignment loss against the Main Branch under a two-stage warmup schedule and a stop-gradient on the index input that confines the auxiliary loss to the index projections. At the 109B-MoE scale, \method{} preserves the capability of a GQA Full-Attention baseline across most pretraining and agentic benchmarks while reducing per-token attention compute by $28.4\times$ at $1\mathrm{M}$ context, the regime in which long-context inference becomes the binding deployment constraint.

\textbf{Outlook.}
\method{}'s core decisions---per-GQA-group independent selection, block-level granularity, and an indexer trained with a KL alignment objective---compose with the GQA backbone shared by most current open-source frontier models, so the recipe should transfer with little modification. Two directions are natural next steps: closing the residual long-context retrieval gap, either through longer sparse training, a larger selection budget at inference time, or a richer indexer scoring function; and extending the same selector-only design to settings beyond pretraining, including reinforcement-learning post-training and agentic deployment, where long-context cost is the dominant operational constraint.
